%% kapa: exam=1 round=24 year=2013 period=2 form=A subject=appraisal_law subject_ko=부동산관계법규 seq=27 number=67
%% source: Q-Net 감정평가사 기출문제 (cst00309) artlSeq=5073802 file=extracted/1008/5073802/2013년 제24회 감정평가사 1차시험 기출문제/2교시(A형).hwp
%% conversion: kordoc 4.24.51 -> tools/gen1.py (auto); review=proofread-sonnet; flags=none
\begin{question}{67}
국유재산법령상 국유재산에 관한 설명으로 옳은 것은?
\begin{choices}
  \item 정부는 정부출자기업체의 고유목적사업을 원활히 수행하기 위하여 자본의 확충이 필요한 경우에는 행정재산과 일반재산을 현물출자할 수 있다.
  \item 행정재산의 관리에 관한 사무를 위임받은 공무원이 경미한 과실로 그 재산에 대하여 손해를 끼친 경우에도 변상의 책임이 있다.
  \item 관리전환하려는 국유재산의 감정평가에 드는 비용이 해당 재산의 가액(價額)에 비하여 과다할 것으로 예상되는 경우 총괄청과 중앙관서의 장 또는 중앙관서의 장 간에 무상으로 관리전환하기로 합의하면 국유재산의 관리전환은 무상으로 할 수 있다.
  \item 국유재산에 관한 사무에 종사하는 직원이 해당 총괄청의 장의 허가 없이 그 처리하는 국유재산을 자기의 소유재산과 교환했다면 총괄청의 장은 이를 취소하여야 한다.
  \item 「국유재산법」에서 정하는 절차와 방법에 따르지 아니하고 일반재산을 사용하거나 수익한 자는 2년 이하의 징역 또는 1천만원 이하의 벌금에 처한다.
\end{choices}
\end{question}
\answer{3}
